\chapter*{Introduction}
\addcontentsline{toc}{chapter}{Introduction}
\markboth{INTRODUCTION}{INTRODUCTION}

\section*{Who this book is for}

This book is written for the working .NET engineer --- someone who has shipped
ASP.NET Core services, knows what a DI container is for, has argued about
where the boundary between application and infrastructure belongs, and is now
being asked to put an agent into production.

It assumes you are comfortable with C\# 12 or later, \code{async}/\code{await},
dependency injection, and Entity Framework Core. It does not assume you have
built an agent before, and it does not assume you have read a paper about
large language models.

It also assumes something less common: that you are sceptical. Most agent
material published since 2023 demonstrates that a thing is \emph{possible}
without ever asking whether it is \emph{reliable}, \emph{observable}, or
\emph{affordable}. A demo that calls three tools and prints a paragraph proves
almost nothing about a system that must run unattended for six hours and hand
work to a human when it is unsure. This book is organised around the second
kind of system.

\section*{What this book is not}

\begin{itemize}[leftmargin=1.4em]
  \item It is not an introduction to large language models. Chapter~\ref{ch:landscape}
        assumes you know roughly what a token, a context window, and a tool call are.
  \item It is not a Semantic Kernel book. Semantic Kernel appears only where it
        explains why an Agent Framework abstraction looks the way it does.
  \item It is not a Python book. The concepts are shared across both runtimes,
        and Python appears twice --- once for cross-runtime interoperability in
        Chapter~\ref{ch:hosting}, once for a feature that has not yet reached .NET
        in Chapter~\ref{ch:harness} --- but every listing that matters is C\#.
  \item It is not prompt engineering. There is exactly one section on writing
        instructions, and it exists because tool descriptions measurably change
        tool-selection accuracy, which is an engineering concern rather than a
        stylistic one.
\end{itemize}

\section*{Why this book exists}

Microsoft Agent Framework reached 1.0 General Availability on 2~April~2026. It is
the convergence of two projects that had accumulated a large amount of field
experience separately: Semantic Kernel's enterprise plumbing and AutoGen's
multi-agent orchestration research. The merged result ships with stable APIs, a
long-term support commitment, and the same programming model across .NET and Python.

That milestone changed the question. Before GA, the interesting question was
\emph{what will this become}. After GA, the interesting question is \emph{how do
you run it}. The documentation answers the first question well and the second
question thinly, because production experience takes time to accumulate and the
framework is only a few months old.

This book is an attempt to write down the second answer while the framework is
new enough that the answer is still worth writing down.

\section*{How the book is organised}

The book has four parts.

\textbf{Part I --- Foundations} (Chapters~\ref{ch:landscape}--\ref{ch:memory})
covers the single-agent surface: where Agent Framework sits in the stack, how an
agent is constructed and run, how tools are bound, and how conversation state and
memory are persisted. If you read only four chapters, read these.

\textbf{Part II --- Workflows} (Chapters~\ref{ch:workflows-basics}--\ref{ch:orchestration})
covers what distinguishes Agent Framework from the majority of agent libraries:
a strongly typed, graph-based, checkpointable execution model, and the
orchestration patterns built on top of it. This is the longest and most
important part of the book.

\textbf{Part III --- Production} (Chapters~\ref{ch:observability}--\ref{ch:harness})
covers middleware, OpenTelemetry, evaluation, hosting, durability, cross-runtime
interoperability, and the agent harness layer added after GA.

\textbf{Part IV --- Capstone} (Chapter~\ref{ch:capstone}) builds one system that
uses the whole of the preceding material, and is deliberately the only chapter
with no new API surface in it.

Chapters are independent enough to read out of order after Part~I, but
Chapters~\ref{ch:workflows-basics} and~\ref{ch:workflows-advanced} must be read
in sequence.

\section*{Conventions}

Four kinds of callout appear throughout.

\begin{note}
Background, context, or a design rationale. Skippable on a first pass.
\end{note}

\begin{warning}
A failure mode that will cost you an afternoon if you meet it unprepared.
Not skippable.
\end{warning}

\begin{versionbox}
Applies to a preview or alpha surface, or to behaviour that is likely to change
between minor versions. Check the release notes before relying on anything
marked this way.
\end{versionbox}

\begin{verifybox}
A listing whose exact signature has been transcribed from documentation or
release notes rather than compiled by the author against the pinned version.
Compile it before you copy it into anything that matters.
\end{verifybox}

Inline, \pkg{Microsoft.Agents.AI} denotes a NuGet package, \api{AsAIAgent} an
API member, and \code{dotnet build} a literal command or file path.

\section*{On versions, and why this book will rot}

Agent Framework does not ship as one package on one version line. It ships as a
family of roughly three dozen packages moving in several \emph{release trains},
and at the time of writing three trains were live simultaneously: the core train
(\pkg{Microsoft.Agents.AI}, \pkg{Microsoft.Agents.AI.OpenAI},
\pkg{Microsoft.Agents.AI.Workflows}) at \mafcore{}; an extensions train
carrying the harness, durable execution and Foundry hosting at \mafext{}; and a
handful of packages stranded on older previews. Packages within a train move
together; packages across trains do not. Assuming one version applies to
everything is the first mistake to avoid.

Every listing in this book is written against the versions recorded on the title
page, and every sample project pins them explicitly through Central Package
Management. Where a feature is newer than GA --- the agent harness, hosted
agents, declarative workflows --- it is marked with a \emph{Version sensitive}
callout.

The honest position is this: the conceptual material in Parts~I and~II will
outlive several minor versions, because agents, tools, executors and edges are
the framework's load-bearing abstractions and Microsoft has committed to their
stability. The wiring details in Part~III will not. Treat the code as a starting
point that you compile, not as scripture.

\section*{The companion repository}

Every listing longer than a few lines exists as a compiling project at:

\begin{center}
\code{github.com/konradcinkusz/maf-book-samples}
\end{center}

The repository is organised by chapter, uses one \code{Directory.Packages.props}
for all pinned versions, and --- with one deliberate exception in
Chapter~\ref{ch:hosting} --- every sample runs against a locally hosted model
with no cloud subscription required. That constraint is not an accessibility
gesture; it is how the samples stay cheap enough to run repeatedly while you are
learning, and it is a constraint the official samples do not meet.

\section*{Environment assumptions}

\begin{itemize}[leftmargin=1.4em]
  \item \dotnetver{} SDK or later.
  \item \vsver{}, JetBrains Rider, or VS~Code with the C\# Dev Kit. Screenshots
        in this book are taken from \vsver{}; the equivalent operations exist in
        the other two.
  \item Docker, for the samples in Chapters~\ref{ch:memory} and~\ref{ch:hosting}.
  \item A locally served model with tool-calling support, for Parts~I and~II.
  \item An Azure subscription with Microsoft Foundry access, for
        Chapters~\ref{ch:hosting} and~\ref{ch:capstone} only.
\end{itemize}

Chapter~\ref{ch:first-agent} walks the setup end to end and does not assume any
of the above is already installed.

\vspace{1.5em}
\noindent\rule{\linewidth}{0.4pt}
\vspace{0.5em}

\noindent\emph{Konrad Cinkusz\\ Mutxamel, \today}

\cleardoublepage
